% Source: https://artofproblemsolving.com/community/c3h1668174_roots_of_an_equation
% Stable identity: aops-topic-1668174-q01
\textbf{Polar form.} Write $z=r(\cos\theta+i\sin\theta)$ with $r>0$. By De Moivre's theorem $z^6=r^6(\cos6\theta+i\sin6\theta)$, and $64=64(\cos0+i\sin0)$. So $r^6=64$, $r=2$, and $6\theta$ is a multiple of $360^\circ$: $\theta=60^\circ k$ for $k=0,1,\dots,5$ (other $k$ repeat these). The six roots are
\[
2(\cos60^\circ k+i\sin60^\circ k):\qquad \boxed{2,\ 1+\sqrt3\,i,\ -1+\sqrt3\,i,\ -2,\ -1-\sqrt3\,i,\ 1-\sqrt3\,i}.
\]
They are the vertices of a regular hexagon inscribed in the circle $|z|=2$:

\begin{asy}
usepackage("amsmath");
size(5.5cm);
defaultpen(fontsize(10pt));
draw((-2.8,0)--(2.8,0),gray,Arrow(TeXHead));draw((0,-2.6)--(0,2.6),gray,Arrow(TeXHead));
label("$\mathrm{Re}$",(2.8,0),E,gray);label("$\mathrm{Im}$",(0,2.6),N,gray);
draw(circle((0,0),2),dashed);
pair[] z;for(int k=0;k<6;++k) z.push(2*dir(60*k));
draw(z[0]--z[1]--z[2]--z[3]--z[4]--z[5]--cycle,blue);
for(int k=0;k<6;++k) dot(z[k],blue);
label("$2$",z[0],SE);label("$1+\sqrt3\,i$",z[1],NE);label("$-1+\sqrt3\,i$",z[2],NW);
label("$-2$",z[3],SW);label("$-1-\sqrt3\,i$",z[4],SW);label("$1-\sqrt3\,i$",z[5],SE);
label("\mbox{}",(min(currentpicture,user=true)+max(currentpicture,user=true))/2,black); // keeps the text after the figure black when LaTeX typesets the labels inline
\end{asy}

\textbf{Check by factoring.} $z^6-64=(z^3-8)(z^3+8)=(z-2)(z^2+2z+4)(z+2)(z^2-2z+4)$. The quadratic factors give $z=-1\pm\sqrt3\,i$ and $z=1\pm\sqrt3\,i$, the same six roots.
